\documentclass[11pt]{article}
\usepackage[T1]{fontenc}
\usepackage{lmodern}
\usepackage{amsmath,amssymb}
\usepackage{microtype}
\usepackage{graphicx}
\usepackage{xcolor}
\usepackage[numbers]{natbib}
\usepackage[colorlinks=true,linkcolor=blue,citecolor=blue!60!black,urlcolor=blue!60!black]{hyperref}

\title{Collective learning from coupled neural parameters}
\author{}
\date{}

\begin{document}
\maketitle

\section{Microscopic dynamics}

Arola-Fernández and Lacasa consider $N$ neural networks, each trained on
private data and coupled through their parameters \cite{arola2024}.
Let $\theta_i^\alpha$ denote one neural parameter of learner $i$ and
$\mathcal L_i$ its local training loss. The update is
\begin{equation}
\theta_i^\alpha(t+1)
=
\theta_i^\alpha(t)
-\eta\,\partial_{\theta_i^\alpha}\mathcal L_i
+\frac{\eta\sigma}{N}
\sum_j q_{ij}\left(\theta_j^\alpha-\theta_i^\alpha\right).
\label{eq:microscopic}
\end{equation}
The private data determine the local gradient, while $q_{ij}$ describes
the interaction network.

\section{Effective theory}

For deep linear networks, Ref.~\cite{arola2024} derives an approximate
reduction to a scalar variable $m_i$ that coarse-grains neural parameters
across layers:
\begin{equation}
\dot m_i
=
\delta_i m_i^D-m_i^{2D+1}-\hat\gamma m_i
+\frac{\hat\sigma}{N}\sum_jq_{ij}(m_j-m_i).
\label{eq:effective}
\end{equation}
Here $D$ is the number of hidden layers,
$\delta_i=\langle \hat x_i\hat y_i\rangle$ is the effective
input-output correlation of the private data, and
$\hat\gamma,\hat\sigma$ are rescaled parameters.
The disorder in $\delta_i$ is quenched, not an adaptive allocation.
The mean effective loss has the scaling
\begin{equation}
\langle\hat L\rangle
\sim \langle m^{2(D+1)}\rangle
-2\langle\delta\rangle\langle m^{D+1}\rangle.
\label{eq:loss}
\end{equation}

\begin{figure}[htb!]
\centering
\includegraphics[width=0.88\linewidth]{figures/fig1.pdf}
\caption{Effective-theory baseline for $D=0,1,2$.
(a) Mean magnetization and (b) normalized effective loss against coupling.
Markers are means over realizations, shaded regions show one standard
deviation, and dash-dotted curves show the adiabatic protocol.}
\label{fig:theory}
\end{figure}

In the neural baseline, each learner trains on a private MNIST digit
class and is evaluated on a common test distribution. The stored neural
runs use a short, nonadiabatic protocol rather than the full published
training horizon.

\begin{figure}[htb!]
\centering
\includegraphics[width=0.88\linewidth]{figures/fig2.pdf}
\caption{Short MNIST baseline averaged over three seeds.
(a) Mean neural parameter and (b) normalized test loss.
These runs do not reproduce the full experimental sweep of
Ref.~\cite{arola2024}.}
\label{fig:mnist}
\end{figure}

\section{Weighted microscopic loss}

To study allocation in a common environment, we must start from a
microscopic loss with sample weights $w_i(x)\geq0$, normalized by
$\langle w_i\rangle=1$:
\begin{equation}
\mathcal L_i^{w}
=
\left\langle w_i(x)\,
\ell\bigl(f(x;\theta_i),y\bigr)\right\rangle
+\gamma\|\theta_i\|^2.
\label{eq:weighted}
\end{equation}
Even for a scalar linear predictor $f(x;u)=ux$ and squared loss,
\begin{align}
\frac12\left\langle w_i(x)(ux-y)^2\right\rangle
&=\frac12 C_i u^2-B_i u+\text{constant},\\
B_i&=\langle w_i(x)xy\rangle,\qquad
C_i=\langle w_i(x)x^2\rangle.
\label{eq:moments}
\end{align}
Thus weighting modifies both the input-output correlation and the
input covariance. The condition $\langle w_i\rangle=1$ does not fix
$C_i$. Consequently, replacing $\delta_i$ in
Eq.~\eqref{eq:effective} with an allocation-dependent coefficient is
not a derived reduction. The weighted deep-linear reduction, including
its assumptions and layerwise dynamics, remains open.

\bibliographystyle{unsrtnat}
\bibliography{refs}
\end{document}
